\documentclass[10pt, twocolumn]{article}
\usepackage{amsmath}
\usepackage{xcolor}
\usepackage[margin=1in]{geometry}
\newcommand{\highlight}[1]{\colorbox{gray}{$\displaystyle#1$}}
\newcommand{\head}[1]{\textnormal{\textbf{#1}}}
\setlength{\columnsep}{.5in}

\begin{document}
	\title{Artificial Neural Networks vs. Support Vector Machines\\for Image Classification}
	\author {Joshua Frankie B. Rayo, Edward Nataniel C. Apostol\\CS 180 THW}
	\maketitle
	
	\section{Summary}
	This paper shows the implementation of image classification using the two most important discriminative models: ANN and SVM.
	Training and testing data are made possible by using ANN and LibSVM libraries, respectively.
	This paper also shows the accuracy of these two discriminative models.
	
	\section{Steps in training and testing data}
	For every subdirectory of 'faces' directory, we removed all the small image files (those filenames that contains \texttt{\_2} or \texttt{\_4} in it).
	The original images are used for training and testing data.
	For the training data, we used only 24 .PGM images per class.
	So how do we form the feature vector? We studied the PGM image specification at \textit{http://netpbm.sourceforge.net/doc/pgm.html}.
	When a PGM image is opened in a plain text editor, we observed the similar pattern:\\
	\texttt{
	P2\\
	24 7\\
	15\\
	0  0  0  0  0  0  0  0  ...\\
	0  3  3  3  3  0  0  7  ...\\
	0  3  0  0  0  0  0  7  ...\\
	0  3  3  3  0  0  0  7  ...\\
	0  3  0  0  0  0  0  7  ...\\
	0  3  0  0  0  0  0  7  ...\\
	0  0  0  0  0  0  0  0  ...\\}
	The first three rows of the image is only a metadata.
	Since we are only concerned about the pixels of the image, we can omit them.
	We form the feature vector row-wise.
	
	\subsection{SVM}
	Given the configuration above, its feature vector to be feed in ANN is
	$F = [<class\_no.> 1:0\: 2:0\: ...\: 25:0\: 26:3\: ...]$.
	The remaining images not trained are used for testing.
	We gathered the images to be trained, obtained an SVM model for it and use the model to test the images to be tested.
	
	\section{Results and Discussions}
	\subsection{SVM}
	For a fixed kernel but different parameters, the accuracy of classification changes.
	Here is the result from different kernels used in SVM:
	\begin{center}
	\begin{tabular}{c c c c}
		\hline
		\head{Kernel} & \head{Param} & \head{Correct} & \head{Accuracy}\\
		\hline
		linear & & 175/180 & 97.2222\% \\
		polynomial & d=3 & 175/180 & 97.2222\% \\
		 & d=6 & 174/180 & 96.6667\% \\
		 & d=10 & 175/180 & 97.2222\% \\
		RBF & g=1/(128*120) & 53/180 & 29.4444\% \\
		 & g=0.1 & 9/180 & 5\% \\
		 & g=0.25 & 9/180 & 5\% \\
		 & g=0.5 & 9/180 & 5\% \\
		 & g=0.75 & 9/180 & 5\% \\
		 & g=1.0 & 9/180 & 5\% \\
		sigmoid & g=1/(128*120) & 9/180 & 5\% \\
		 & g=0.1 & 9/180 & 5\% \\
		 & g=0.25 & 9/180 & 5\% \\
		 & g=0.5 & 9/180 & 5\% \\
		 & g=0.75 & 9/180 & 5\% \\
		 & g=1.0 & 9/180 & 5\% \\
		\hline
	\end{tabular}
	\end{center}
	It is a very good idea to treat each pixel as a feature, because each person has different facial features.
	It is very important to notice the person's difference with each other.
	Another good thing is, even if that person uses sunglasses or change his facial expression, the SVM can nearly recognize the face of that person.
	For the sigmoid kernel, which appear to be the worst kernel, it does not seem to have a change.
	We can say that the linear function is the best SVM kernel for image classification (for PGM files).
\end{document}